% --- Archon physics formalization source begin ---
% archon:physics
% archon:covers PhyXMiniProblems/problem_phyx_mini_0182.lean
% archon:source-report reports/phyx_mini/problem_phyx_mini_0182.source.json

\chapter{Physics problem phyx\_mini\_0182}
\label{ch:PhyXMiniProblems_problem_phyx_mini_0182}

\paragraph{Problem source.}
\#\# Physical scenario

The two strings in figure are of equal length and are being driven at equal frequencies. The linear density of the left string is \$\textbackslash{}mu\_0\$.

\#\# Question

What is the linear density of the right string?

\#\# Answer choices

- A: A: \$\textbackslash{}mu\_0 \textbackslash{}frac\{4\}\{9\}\$

- B: B: \$\textbackslash{}mu\_0 \textbackslash{}frac\{3\}\{2\}\$

- C: C: \$\textbackslash{}mu\_0 \textbackslash{}frac\{2\}\{3\}\$

- D: D: \$\textbackslash{}mu\_0 \textbackslash{}frac\{9\}\{4\}\$

\#\# Figure caption (auxiliary; use the image as primary evidence)

The image is a diagram featuring a stretched spring in the center, labeled "Stretched spring." On either side of the spring, sinusoidal wave patterns are depicted, suggesting wave transmission through the spring. The waves are symmetric and mirror each other. The ends of the waves are connected to vertical barriers or surfaces, indicating fixed endpoints. The overall setup suggests a physical representation of wave motion in a medium with a spring-like property.

\#\# Dataset metadata

Wave Properties; Physical Model Grounding Reasoning, Multi-Formula Reasoning

\paragraph{Recorded answer/context.}
D: D: \$\textbackslash{}mu\_0 \textbackslash{}frac\{9\}\{4\}\$

\paragraph{Figure/image path.}
/root/proposal\_for\_physic/science-mango/phyx\_mini\_run/phyx\_data/test\_image/182.png

\paragraph{Formalization target.}
create a compiling Lean file with sorry bodies at `PhyXMiniProblems/problem\_phyx\_mini\_0182.lean`. The Lean declarations must preserve the physical quantities, dimensions or dimensional roles, figure labels, governing-law hypotheses, and final relation expressed by this problem.
Use Mathlib/Physlib names found through LeanExplore where available. If a physics API is missing, introduce faithful local abstractions rather than scalar placeholder aliases.

\begin{theorem}[Physics formalization target]
\label{thm:physics:phyx_mini_0182:target}
\lean{PhyXMiniProblems.ProblemPhyXMini0182.problem_phyx_mini_0182}
\uses{def:physics:phyx-mini-0182:phyxminiproblems-problemphyxmini0182-linearmassdensity, def:physics:phyx-mini-0182:phyxminiproblems-problemphyxmini0182-coupledstringdiagram, def:physics:phyx-mini-0182:phyxminiproblems-problemphyxmini0182-matchesproblemstatement, def:physics:phyx-mini-0182:phyxminiproblems-problemphyxmini0182-matchessuppliedfigure, def:physics:phyx-mini-0182:phyxminiproblems-problemphyxmini0182-hasphysicalstringparameters, def:physics:phyx-mini-0182:phyxminiproblems-problemphyxmini0182-satisfiescoupledstringstandingwavelaws}
The assigned autoformalize agent should translate this physics problem into Lean declarations in the covered file, with theorem and lemma proof bodies written as `by sorry`.
\end{theorem}
\begin{proof}
This is an autoformalization task, not a proof task. Produce statements that can later be proved by the physics prover without weakening the physical model.
\end{proof}

% --- Archon declaration topology begin ---
\section{Declaration topology}
\begin{definition}[Length Quantity]
  \label{def:physics:phyx-mini-0182:phyxminiproblems-problemphyxmini0182-lengthquantity}
  \lean{PhyXMiniProblems.ProblemPhyXMini0182.LengthQuantity}
  Physical length, used for both string length and standing wavelength
\end{definition}
\begin{proof}
  This is the declared model definition of Length Quantity.
\end{proof}
\begin{definition}[Linear Mass Density]
  \label{def:physics:phyx-mini-0182:phyxminiproblems-problemphyxmini0182-linearmassdensity}
  \lean{PhyXMiniProblems.ProblemPhyXMini0182.LinearMassDensity}
  Linear mass density, with dimension mass per length
\end{definition}
\begin{proof}
  This is the declared model definition of Linear Mass Density.
\end{proof}
\begin{definition}[Frequency Quantity]
  \label{def:physics:phyx-mini-0182:phyxminiproblems-problemphyxmini0182-frequencyquantity}
  \lean{PhyXMiniProblems.ProblemPhyXMini0182.FrequencyQuantity}
  Ordinary frequency, with dimension inverse time
\end{definition}
\begin{proof}
  This is the declared model definition of Frequency Quantity.
\end{proof}
\begin{definition}[Speed Quantity]
  \label{def:physics:phyx-mini-0182:phyxminiproblems-problemphyxmini0182-speedquantity}
  \lean{PhyXMiniProblems.ProblemPhyXMini0182.SpeedQuantity}
  Transverse-wave propagation speed, using Physlib's unit-aware speed type
\end{definition}
\begin{proof}
  This is the declared model definition of Speed Quantity.
\end{proof}
\begin{definition}[Tension Quantity]
  \label{def:physics:phyx-mini-0182:phyxminiproblems-problemphyxmini0182-tensionquantity}
  \lean{PhyXMiniProblems.ProblemPhyXMini0182.TensionQuantity}
  String tension, with the physical dimension of force
\end{definition}
\begin{proof}
  This is the declared model definition of Tension Quantity.
\end{proof}
\begin{definition}[length Readout]
  \label{def:physics:phyx-mini-0182:phyxminiproblems-problemphyxmini0182-lengthreadout}
  \lean{PhyXMiniProblems.ProblemPhyXMini0182.lengthReadout}
  \uses{def:physics:phyx-mini-0182:phyxminiproblems-problemphyxmini0182-lengthquantity}
  Scalar readout of a physical length in a coherent choice of units
\end{definition}
\begin{proof}
  This is the declared model definition of length Readout.
\end{proof}
\begin{definition}[linear Density Readout]
  \label{def:physics:phyx-mini-0182:phyxminiproblems-problemphyxmini0182-lineardensityreadout}
  \lean{PhyXMiniProblems.ProblemPhyXMini0182.linearDensityReadout}
  \uses{def:physics:phyx-mini-0182:phyxminiproblems-problemphyxmini0182-linearmassdensity}
  Scalar readout of a linear mass density in a coherent choice of units
\end{definition}
\begin{proof}
  This is the declared model definition of linear Density Readout.
\end{proof}
\begin{definition}[frequency Readout]
  \label{def:physics:phyx-mini-0182:phyxminiproblems-problemphyxmini0182-frequencyreadout}
  \lean{PhyXMiniProblems.ProblemPhyXMini0182.frequencyReadout}
  \uses{def:physics:phyx-mini-0182:phyxminiproblems-problemphyxmini0182-frequencyquantity}
  Scalar readout of an ordinary frequency in a coherent choice of units
\end{definition}
\begin{proof}
  This is the declared model definition of frequency Readout.
\end{proof}
\begin{definition}[speed Readout]
  \label{def:physics:phyx-mini-0182:phyxminiproblems-problemphyxmini0182-speedreadout}
  \lean{PhyXMiniProblems.ProblemPhyXMini0182.speedReadout}
  \uses{def:physics:phyx-mini-0182:phyxminiproblems-problemphyxmini0182-speedquantity}
  Scalar readout of a wave speed in a coherent choice of units
\end{definition}
\begin{proof}
  This is the declared model definition of speed Readout.
\end{proof}
\begin{definition}[tension Readout]
  \label{def:physics:phyx-mini-0182:phyxminiproblems-problemphyxmini0182-tensionreadout}
  \lean{PhyXMiniProblems.ProblemPhyXMini0182.tensionReadout}
  \uses{def:physics:phyx-mini-0182:phyxminiproblems-problemphyxmini0182-tensionquantity}
  Scalar readout of a string tension in a coherent choice of units
\end{definition}
\begin{proof}
  This is the declared model definition of tension Readout.
\end{proof}
\begin{definition}[String Side]
  \label{def:physics:phyx-mini-0182:phyxminiproblems-problemphyxmini0182-stringside}
  \lean{PhyXMiniProblems.ProblemPhyXMini0182.StringSide}
  The two strings on opposite sides of the central spring
\end{definition}
\begin{proof}
  This is the declared model definition of String Side.
\end{proof}
\begin{definition}[String End]
  \label{def:physics:phyx-mini-0182:phyxminiproblems-problemphyxmini0182-stringend}
  \lean{PhyXMiniProblems.ProblemPhyXMini0182.StringEnd}
  The end of a string nearest the wall or nearest the central connector
\end{definition}
\begin{proof}
  This is the declared model definition of String End.
\end{proof}
\begin{definition}[Endpoint Attachment]
  \label{def:physics:phyx-mini-0182:phyxminiproblems-problemphyxmini0182-endpointattachment}
  \lean{PhyXMiniProblems.ProblemPhyXMini0182.EndpointAttachment}
  Mechanical attachment shown at a string endpoint
\end{definition}
\begin{proof}
  This is the declared model definition of Endpoint Attachment.
\end{proof}
\begin{definition}[Transverse Boundary Condition]
  \label{def:physics:phyx-mini-0182:phyxminiproblems-problemphyxmini0182-transverseboundarycondition}
  \lean{PhyXMiniProblems.ProblemPhyXMini0182.TransverseBoundaryCondition}
  Transverse-displacement condition visible at a standing-wave endpoint
\end{definition}
\begin{proof}
  This is the declared model definition of Transverse Boundary Condition.
\end{proof}
\begin{definition}[Central Connector Kind]
  \label{def:physics:phyx-mini-0182:phyxminiproblems-problemphyxmini0182-centralconnectorkind}
  \lean{PhyXMiniProblems.ProblemPhyXMini0182.CentralConnectorKind}
  The connector explicitly labeled in the center of the supplied figure
\end{definition}
\begin{proof}
  This is the declared model definition of Central Connector Kind.
\end{proof}
\begin{definition}[Coupled String Diagram]
  \label{def:physics:phyx-mini-0182:phyxminiproblems-problemphyxmini0182-coupledstringdiagram}
  \lean{PhyXMiniProblems.ProblemPhyXMini0182.CoupledStringDiagram}
  \uses{def:physics:phyx-mini-0182:phyxminiproblems-problemphyxmini0182-lengthquantity, def:physics:phyx-mini-0182:phyxminiproblems-problemphyxmini0182-linearmassdensity, def:physics:phyx-mini-0182:phyxminiproblems-problemphyxmini0182-frequencyquantity, def:physics:phyx-mini-0182:phyxminiproblems-problemphyxmini0182-speedquantity, def:physics:phyx-mini-0182:phyxminiproblems-problemphyxmini0182-tensionquantity, def:physics:phyx-mini-0182:phyxminiproblems-problemphyxmini0182-stringside, def:physics:phyx-mini-0182:phyxminiproblems-problemphyxmini0182-stringend, def:physics:phyx-mini-0182:phyxminiproblems-problemphyxmini0182-endpointattachment, def:physics:phyx-mini-0182:phyxminiproblems-problemphyxmini0182-transverseboundarycondition, def:physics:phyx-mini-0182:phyxminiproblems-problemphyxmini0182-centralconnectorkind}
  All physical quantities and figure labels for the coupled-string system. halfWavelengthCount is the number of loops (antinodes) visible along a string. The right linear density is an unconstrained physical quantity in this structure; no answer value is built into the setup
\end{definition}
\begin{proof}
  This is the declared model definition of Coupled String Diagram.
\end{proof}
\begin{definition}[Matches Problem Statement]
  \label{def:physics:phyx-mini-0182:phyxminiproblems-problemphyxmini0182-matchesproblemstatement}
  \lean{PhyXMiniProblems.ProblemPhyXMini0182.MatchesProblemStatement}
  \uses{def:physics:phyx-mini-0182:phyxminiproblems-problemphyxmini0182-linearmassdensity, def:physics:phyx-mini-0182:phyxminiproblems-problemphyxmini0182-coupledstringdiagram}
  The data stated in the prose: equal physical lengths, equal drive frequencies, and left-string linear density μ₀. No relation for the right density occurs here
\end{definition}
\begin{proof}
  This is the declared model definition of Matches Problem Statement.
\end{proof}
\begin{definition}[Matches Supplied Figure]
  \label{def:physics:phyx-mini-0182:phyxminiproblems-problemphyxmini0182-matchessuppliedfigure}
  \lean{PhyXMiniProblems.ProblemPhyXMini0182.MatchesSuppliedFigure}
  \uses{def:physics:phyx-mini-0182:phyxminiproblems-problemphyxmini0182-coupledstringdiagram}
  Primary-image readout: two loops appear on the left, three on the right, both outer ends meet fixed walls, both inner ends meet the stretched spring, and all four depicted standing-wave endpoints are displacement nodes
\end{definition}
\begin{proof}
  This is the declared model definition of Matches Supplied Figure.
\end{proof}
\begin{definition}[Has Physical String Parameters]
  \label{def:physics:phyx-mini-0182:phyxminiproblems-problemphyxmini0182-hasphysicalstringparameters}
  \lean{PhyXMiniProblems.ProblemPhyXMini0182.HasPhysicalStringParameters}
  \uses{def:physics:phyx-mini-0182:phyxminiproblems-problemphyxmini0182-lengthreadout, def:physics:phyx-mini-0182:phyxminiproblems-problemphyxmini0182-lineardensityreadout, def:physics:phyx-mini-0182:phyxminiproblems-problemphyxmini0182-frequencyreadout, def:physics:phyx-mini-0182:phyxminiproblems-problemphyxmini0182-speedreadout, def:physics:phyx-mini-0182:phyxminiproblems-problemphyxmini0182-tensionreadout, def:physics:phyx-mini-0182:phyxminiproblems-problemphyxmini0182-coupledstringdiagram}
  Positivity conditions selecting the ordinary nondegenerate string regime
\end{definition}
\begin{proof}
  This is the declared model definition of Has Physical String Parameters.
\end{proof}
\begin{definition}[Satisfies Coupled String Standing Wave Laws]
  \label{def:physics:phyx-mini-0182:phyxminiproblems-problemphyxmini0182-satisfiescoupledstringstandingwavelaws}
  \lean{PhyXMiniProblems.ProblemPhyXMini0182.SatisfiesCoupledStringStandingWaveLaws}
  \uses{def:physics:phyx-mini-0182:phyxminiproblems-problemphyxmini0182-lengthreadout, def:physics:phyx-mini-0182:phyxminiproblems-problemphyxmini0182-lineardensityreadout, def:physics:phyx-mini-0182:phyxminiproblems-problemphyxmini0182-frequencyreadout, def:physics:phyx-mini-0182:phyxminiproblems-problemphyxmini0182-speedreadout, def:physics:phyx-mini-0182:phyxminiproblems-problemphyxmini0182-tensionreadout, def:physics:phyx-mini-0182:phyxminiproblems-problemphyxmini0182-coupledstringdiagram}
  Standard laws used for the two strings, expressed in every coherent unit choice: a node-to-node string containing n loops has 2 L = n λ; wave kinematics gives v = f λ; a stretched string obeys T = μ v²; the same massless stretched spring transmits equal tension magnitudes to its two ends. These are general physical laws. None mentions the problem-specific density factor 9/4 or any answer choice
\end{definition}
\begin{proof}
  This is the declared model definition of Satisfies Coupled String Standing Wave Laws.
\end{proof}
\begin{lemma}[linear Density harmonic ratio]
  \label{lem:physics:phyx-mini-0182:phyxminiproblems-problemphyxmini0182-lineardensity-harmonic-ratio}
  \lean{PhyXMiniProblems.ProblemPhyXMini0182.linearDensity_harmonic_ratio}
  \uses{def:physics:phyx-mini-0182:phyxminiproblems-problemphyxmini0182-linearmassdensity, def:physics:phyx-mini-0182:phyxminiproblems-problemphyxmini0182-lineardensityreadout, def:physics:phyx-mini-0182:phyxminiproblems-problemphyxmini0182-coupledstringdiagram, def:physics:phyx-mini-0182:phyxminiproblems-problemphyxmini0182-matchesproblemstatement, def:physics:phyx-mini-0182:phyxminiproblems-problemphyxmini0182-hasphysicalstringparameters, def:physics:phyx-mini-0182:phyxminiproblems-problemphyxmini0182-satisfiescoupledstringstandingwavelaws}
  The generic consequence needed before substituting the figure's loop counts: for equal length, frequency, and tension, density times the square of the other string's relevant harmonic denominator gives the cross-multiplied ratio
\end{lemma}
\begin{proof}
  The conclusion follows from the governing relations and figure readouts named in the statement of linear Density harmonic ratio.
\end{proof}
\begin{definition}[Answer Choice]
  \label{def:physics:phyx-mini-0182:phyxminiproblems-problemphyxmini0182-answerchoice}
  \lean{PhyXMiniProblems.ProblemPhyXMini0182.AnswerChoice}
  Labels of the four answer choices printed in the problem
\end{definition}
\begin{proof}
  This is the declared model definition of Answer Choice.
\end{proof}
\begin{definition}[displayed Density Multiplier]
  \label{def:physics:phyx-mini-0182:phyxminiproblems-problemphyxmini0182-displayeddensitymultiplier}
  \lean{PhyXMiniProblems.ProblemPhyXMini0182.displayedDensityMultiplier}
  \uses{def:physics:phyx-mini-0182:phyxminiproblems-problemphyxmini0182-answerchoice}
  Dimensionless multiplier of μ₀ printed beside each answer choice
\end{definition}
\begin{proof}
  This is the declared model definition of displayed Density Multiplier.
\end{proof}
\begin{definition}[displayed Linear Density]
  \label{def:physics:phyx-mini-0182:phyxminiproblems-problemphyxmini0182-displayedlineardensity}
  \lean{PhyXMiniProblems.ProblemPhyXMini0182.displayedLinearDensity}
  \uses{def:physics:phyx-mini-0182:phyxminiproblems-problemphyxmini0182-linearmassdensity, def:physics:phyx-mini-0182:phyxminiproblems-problemphyxmini0182-answerchoice, def:physics:phyx-mini-0182:phyxminiproblems-problemphyxmini0182-displayeddensitymultiplier}
  The physical density displayed by an answer choice
\end{definition}
\begin{proof}
  This is the declared model definition of displayed Linear Density.
\end{proof}
% --- Archon declaration topology end ---
% --- Archon physics formalization source end ---
